\subsection{THE AUTOMORPHISM GROUP OF A LIE GROUP (REAL OR COMPLEX CASE)}

In this no., we assume that K = R or C.

Lemma 4. Let \(\mathbf{H}\) be a finite-dimensional simply connected Lie group.

\begin{enumerate}
\def\labelenumi{(\roman{enumi})}
\item
  For all \(u \in \operatorname{Aut} \mathbf{L}(\mathbf{H})\) , let \(\theta(u)\) be the unique automorphism of \(\mathbf{H}\) such that \(\mathbf{L}(\theta(u)) = u\) . Then the mapping \((u, g) \mapsto \theta(u)g\) of \((\operatorname{Aut} \mathbf{L}(\mathbf{H})) \times \mathbf{H}\) into \(\mathbf{H}\) is analytic.
\item
  Let \(\mathbf{N}\) be a Lie subgroup of \(\mathbf{H}\) and \(\operatorname{Aut}(\mathbf{H},\mathbf{N})\) the set of \(v\in \operatorname{Aut}\mathbf{H}\) such that \(v(\mathbf{N}) = \mathbf{N}\) . Then \(\theta^{-1}(\operatorname {Aut}(\mathbf{H},\mathbf{N}))\) is a Lie subgroup of \(\operatorname {Aut}\mathbf{L}(\mathbf{H})\) .
\item
  Suppose that N is discrete and normal, so that the Lie algebra of G = H/N is identified with L(H). For all \(w \in Aut G\) , let \(\eta(w)\) be the unique automorphism of H such that \(\mathrm{L}(\eta(w)) = \mathrm{L}(w)\) . Then the mapping \(\eta\) is an isomorphism of the group Aut G onto the group \(\mathrm{Aut}(\mathrm{H}, \mathrm{N})\) .
\end{enumerate}

To prove (i), it suffices, by Lemma 3 of no. 1, to verify that the mapping \((u,g)\mapsto\theta(u)g\) is analytic in a neighbourhood of \((\mathrm{Id}_{\mathrm{L(H)}},e)\) . There exists an open neighbourhood B of 0 in L(H) such that \(\psi = \exp_{\mathbf{H}}|\mathbf{B}\) is an analytic isomorphism of B onto an open neighbourhood of \(e\) in H. There exist an open neighbourhood U of \(\mathrm{Id}_{\mathrm{L(H)}}\) in Aut L(H) and an open neighbourhood B' of 0 in L(H) such that \(\mathrm{U(B')} \subset \mathrm{B}\) . Then the mapping \((u, g) \mapsto \theta(u)g\) of \(\mathrm{U} \times \psi(\mathrm{B'})\) into H is composed of the following mappings:

\[
\begin{array}{l l l l} \text {the mapping} (u, g) \mapsto \langle u, \psi^ {- 1} (g) \rangle & \text {of} & \mathrm{U} \times \psi (\mathrm{B} ^ {\prime}) & \text {into} \quad \mathrm{U} \times \mathrm{B} ^ {\prime}; \\ \text {the mapping} (u, x) \mapsto u (x) & \text {of} & \mathrm{U} \times \mathrm{B} ^ {\prime} & \text {into} \quad \mathrm{B}; \\ \text {the mapping} y \mapsto \psi (y) & \text {of} & \mathrm{B} & \text {into} \quad \mathrm{G}. \end{array}
\]

Hence this mapping is analytic.

Let \(p\) be the canonical mapping of \(H\) into the homogeneous space \(H/N\) . Then \(\theta^{-1}(\text{Aut}(H, N))\) is the set of \(u \in \text{Aut } L(H)\) such that

\[
p (\theta (u) g) = p (e), \quad p (\theta (u ^ {- 1}) g) = p (e)
\]

for all \(g \in \mathbb{N}\) . By § 8, no. 2, Theorem 2 and Corollary 2 to Theorem 2, this proves (ii).

Suppose that N is discrete and normal. Let \(w \in \operatorname{Aut} G\) . Then

\[
\mathrm{L} (p \circ \eta (w)) = \mathrm{L} (\eta (w)) = \mathrm{L} (w) = \mathrm{L} (w \circ p)
\]

hence \(p \circ \eta(w) = w \circ p\) and therefore \(\eta(w) \in \text{Aut(H, N)}\) . Clearly the mapping \(\eta\) of Aut G into Aut(H, N) is an injective homomorphism. This homomorphism is surjective because \(p: H \to G\) is a submersion.

Let G be a locally compact group and \(\Gamma\) the automorphism group of G. Recall that a topology \(T_{\beta}\) has been defined on \(\Gamma\) (General Topology, Chapter X, §3, no. 5). It is the coarsest topology for which the mappings \(v \mapsto v\) and \(v \mapsto v^{-1}\) of \(\Gamma\) into \(\mathcal{C}_{\varepsilon}(G; G)\) (space of continuous mappings of G into G with the compact convergence topology) are continuous. The topology \(T_{\beta}\) is compatible with the group structure on \(\Gamma\) (loc. cit.). For every compact subset L of G and every neighbourhood U of \(e_{\mathrm{G}}\) in G, let N(L, U) be the set of \(\phi \in \Gamma\) such that \(\phi(g) \in gU\) and \(\phi^{-1}(g) \in gU\) for all \(g \in L\) ; then the N(L, U) form a fundamental system of neighbourhoods of \(e_{\Gamma}\) . If G is generated by a compact subset C, the topology \(T_{\beta}\) is also the coarsest topology for which the mappings \(v \mapsto v[C\) and \(v \mapsto v^{-1}[C]\) of \(\Gamma\) into \(\mathcal{C}_{u}(C; G)\) are continuous (for every compact subset of G is contained in \((\mathrm{C} \cup \mathrm{C}^{-1})^{\mathrm{n}}\) for sufficiently large n). If K is locally compact and V is a finite-dimensional vector space over K, the topology \(T_{\beta}\) on GL(V) is just the usual topology.

THEOREM 1. Let G be a finite-dimensional Lie group and \(G_{0}\) its identity component. Suppose that G is generated by \(G_{0}\) and a finite number of elements.

\begin{enumerate}
\def\labelenumi{(\roman{enumi})}
\tightlist
\item
  There exists on Aut G one and only one analytic manifold structure satisfying the following condition:
\end{enumerate}

(AUT) for every analytic manifold M and every mapping f of M into Aut G, f is analytic if and only if the mapping \((m, g) \mapsto f(m)g\) of M × G into G is analytic.

Suppose in the rest of the statement that Aut G has this structure.

\begin{enumerate}
\def\labelenumi{(\roman{enumi})}
\setcounter{enumi}{1}
\item
  Aut G is a finite-dimensional Lie group.
\item
  The morphism \(\phi: u \mapsto L(u)\) of Aut G into Aut L(G) is analytic.
\item
  If G is connected, \(\phi\) is an isomorphism of the Lie group Aut G onto a Lie subgroup of Aut L(G); this Lie subgroup is equal to Aut L(G) if G is simply connected.
\item
  Let \(\mathfrak{a}\) be the set of infinitesimal automorphisms of \(\mathbf{G}\) . Then \(\mathfrak{a}\) is a Lie algebra of vector fields and the law of infinitesimal operation associated with the mapping \((u, g) \mapsto u(g)\) of (Aut G) \(\times\) G into G is an isomorphism of L(Aut G) onto \(\mathfrak{a}\) .
\item
  The topology of the Lie group Aut G is the topology \(\mathcal{T}_{\beta}\) .
\end{enumerate}

\begin{enumerate}
\def\labelenumi{(\alph{enumi})}
\item
  The uniqueness of the analytic structure considered in (i) is obvious.
\item
  Suppose that G is connected. Let H be the universal covering space of G, p the canonical morphism of H onto G and N = Ker p.~We introduce the notation \(\theta\) , \(\eta\) and \(\operatorname{Aut}(\mathrm{H},\mathrm{N})\) of Lemma 4. We transport the Lie group structure of Aut L(G) to Aut H by means of \(\theta\) . Then Aut H becomes a finite-dimensional Lie group and \(\operatorname{Aut}(\mathrm{H},\mathrm{N})\) a Lie subgroup of Aut H (Lemma 4 (ii)). We transport the Lie group structure of \(\operatorname{Aut}(\mathrm{H},\mathrm{N})\) to Aut G by means of \(\eta^{-1}\) . Then Aut G becomes a finite-dimensional Lie group. Properties (ii), (iii) and (iv) of the theorem are satisfied and the mapping \((u,g)\mapsto u(g)\) of \((\operatorname{Aut}\mathbf{G})\times\mathbf{G}\) into G is analytic (Lemma 4 (i)). Let M be an analytic manifold, f a mapping of M into Aut G and \(\phi\) the mapping \((m,g)\mapsto f(m)g\) of M × G into G. Clearly, if f is analytic, \(\phi\) is analytic. Suppose that \(\phi\) is analytic. Then the T \(\phi\) : TM × TG → TG is analytic; its restriction to M × L(G), that is the mapping \((m,x)\mapsto\operatorname{L}(f(m))x\) of M × L(G) into L(G) is therefore analytic; as L(G) is finite-dimensional, it follows that the mapping \(m\mapsto\operatorname{L}(f(m))\) of M into Aut L(G) is analytic and hence that f is analytic. Thus (i) holds.
\end{enumerate}

Let \(\mathrm{L}(\mathrm{G})\) be given a norm. For all \(\lambda > 0\) , let \(\mathrm{B}_{\lambda}\) be the open ball of centre 0 and radius \(\lambda\) in \(\mathrm{L}(\mathrm{G})\) . We choose \(\lambda > 0\) sufficiently small for \(\psi = \exp_{\mathrm{G}}|\mathrm{B}_{\lambda}\) to be an isomorphism of the analytic manifold \(\mathrm{B}_{\lambda}\) onto the open submanifold \(\psi(\mathrm{B}_{\lambda})\) of \(G\) . Let \(\Phi\) be a filter on Aut \(G\) . For \(\Phi\) to converge to \(\mathrm{Id}_{\mathrm{G}}\) in Aut \(G\) , it is necessary and sufficient that \(\mathrm{L}(\Phi)\) converge to \(\mathrm{Id}_{\mathrm{L}(\mathrm{G})}\) in Aut \(L(\mathrm{G})\) and hence that \(L(\Phi)|B_{\lambda/2}\) and \(L(\Phi)^{-1}\mathrm{B}_{\lambda/2}\) converge uniformly to \(\mathrm{Id}_{\mathrm{B}_{\lambda/2}}\) . This condition implies that \(\Phi|\psi(\mathrm{B}_{\lambda/2})\) and \(\Phi^{-1}|\psi(\mathrm{B}_{\lambda/2})\) converge uniformly to \(\mathrm{Id}_{\psi(\mathrm{B}_{\lambda/2})}\) . Conversely, suppose that \(\Phi|\psi(\mathrm{B}_{\lambda/2})\) converges uniformly to \(\mathrm{Id}_{\psi(\mathrm{B}_{\lambda/2})}\) . There exists \(M \in \Phi\) such that, if \(u \in M\) , then \(u(\psi(\mathrm{B}_{\lambda/2})) \subset \psi(\mathrm{B}_{2\lambda/3})\) ; then \(L(u)(\mathrm{B}_{\lambda/2})\) is a connected subset of \(L(\mathrm{G})\) whose image under \(\exp_{\mathrm{G}}\) is contained in \(\psi(\mathrm{B}_{2\lambda/3})\) , hence \(L(u)(\mathrm{B}_{\lambda/2})\) does not meet \(B_{\lambda} - B_{2\lambda/3}\) and therefore \(L(u)(\mathrm{B}_{\lambda/2}) \subset B_{\lambda}\) ; then the hypothesis that \(\Phi|\psi(\mathrm{B}_{\lambda/2})\) converges uniformly to \(\mathrm{Id}_{\psi(\mathrm{B}_{\lambda/2})}\) implies that \(L(\Phi)|B_{\lambda/2}\) converges uniformly to \(\mathrm{Id}_{\mathrm{B}_{\lambda/2}}\) . It then follows that:

\((\Phi\) converges to \(\mathrm{Id}_{\mathrm{G}}\) in Aut G) \(\Leftrightarrow\) ( \(\Phi\) converges to \(\mathrm{Id}_{\mathrm{G}}\) under \(\mathcal{T}_{\beta}\) ).

This proves (vi).

Let D be the law of infinitesimal operation associated with the law of left operation of \(\mathrm{Aut(G)}\) on G. By Propositions 1 and 2 of no. 1, \(\mathrm{D(L(AutG))} = a\) . Hence \(a\) is a Lie algebra of vector fields and D is a morphism of \(\mathrm{L(AutG)}\) onto \(a\) . Let \(x_{1}\) and \(x_{2}\) be elements of \(\mathrm{L(AutG)}\) such that \(\mathrm{D}(x_1) = \mathrm{D}(x_2)\) . Then the laws of operation \((\lambda, g) \mapsto (\exp \lambda x_1)g\) and \((\lambda, g) \mapsto (\exp \lambda x_2)g\) of K on G have the same associated law of infinitesimal operation; hence, for \(|\lambda|\) sufficiently small, \(\exp \lambda x_1\) and \(\exp \lambda x_2\) coincide on a neighbourhood of \(e\) (§ 4, no. 7, Theorem 6), whence \(\exp \lambda x_1 = \exp \lambda x_2\) . It follows that \(x_1 = x_2\) and hence D is an isomorphism of \(\mathrm{L(AutG)}\) onto \(a\) .

The theorem has thus been completely proved for G connected.

\begin{enumerate}
\def\labelenumi{(\alph{enumi})}
\setcounter{enumi}{2}
\tightlist
\item
  We pass to the general case. By hypothesis, G is generated by \(G_{0}\) and a finite number of elements \(x_{1}, x_{2}, \ldots, x_{n}\) . Every \(u \in AutG\) leaves \(G_{0}\) stable. Let \(Aut_{1}G\) be the set of \(u \in AutG\) which, on passing to the quotient, give the identity automorphism of \(G/G_{0}\) . This is a normal subgroup of Aut G. By part (b) of the proof, \(AutG_{0}\) has a canonical Lie group structure and the mapping \((g_{1}, g_{2}, \ldots, g_{n}, u) \mapsto (ug_{1}, ug_{2}, \ldots, ug_{n})\) of \(G_{0}^{n} \times AutG_{0}\) into \(G_{0}^{n}\) is analytic. Let P be the corresponding semidirect product of \(AutG_{0}\) by \(G_{0}^{n}\) ; it is a finite-dimensional Lie group ( \(\S\) 1, no. 4, Proposition 7).
\end{enumerate}

If \(w\in \mathrm{Aut}_1\mathrm{G}\) , we write

\[
\begin{array}{r l} w _ {0} & = w | \mathrm{G} _ {0} \in \mathrm{Aut}   \mathrm{G} _ {0} \\ w _ {i} & = x _ {i} ^ {- 1} w (x _ {i}) \in \mathrm{G} _ {0} \quad (1 \leqslant i \leqslant n) \\ \zeta (w) & = ((w _ {1}, \ldots , w _ {n}), w _ {0}) \in \mathrm{P}. \end{array}
\]

For all \(w, w'\) in \(\mathrm{Aut}_1\mathrm{G}\) ,

\[
\begin{array}{r l} \zeta (w) \zeta (w ^ {\prime}) & = \left((w _ {1}, \ldots , w _ {n}), w _ {0}\right) \left((w _ {1} ^ {\prime}, \ldots , w _ {n} ^ {\prime}), w _ {0} ^ {\prime}\right) \\ & = \left((w _ {1} w _ {0} (w _ {1} ^ {\prime}), \ldots , w _ {n} w _ {0} (w _ {n} ^ {\prime})), w _ {0} w _ {0} ^ {\prime}\right) \\ & = \left((x _ {1} ^ {- 1} w (x _ {1}) w (x _ {1} ^ {- 1} w ^ {\prime} (x _ {1})), \ldots , x _ {n} ^ {- 1} w (x _ {n}) w (x _ {n} ^ {- 1} w ^ {\prime} (x _ {n}))), w _ {0} w _ {0} ^ {\prime}\right) \\ & = \left(((w w ^ {\prime}) _ {1}, \ldots , (w w ^ {\prime}) _ {n}), (w w ^ {\prime}) _ {0}\right) \\ & = \zeta (w w ^ {\prime}) \end{array}
\]

and hence \(\zeta\) is a homomorphism of \(\mathrm{Aut}_1\mathrm{G}\) into P. This homomorphism is obviously injective.

We show that \(\zeta(\mathrm{Aut}_1\mathrm{G})\) is closed in P. Let \(\Phi\) be a filter on \(\mathrm{Aut}_1\mathrm{G}\) such that \(\zeta(\Phi)\) converges to a point \(((w_1,\dots,w_n),w_0)\) of P. Then \(\Phi\) converges pointwise to a mapping \(v\) of G into G. Clearly \(v\) is an endomorphism of the group G. Moreover, \(v\) leaves each coset modulo \(\mathbf{G}_0\) stable and \(v |\mathrm{G}_0 = w_0\) . It follows that \(v\in \mathrm{Aut}_1\mathrm{G}\) . As \(\zeta(v) = ((w_1,\dots,w_n),w_0)\) , we have shown that \(\zeta(\mathrm{Aut}_1\mathrm{G})\) is closed in P.

\begin{enumerate}
\def\labelenumi{(\alph{enumi})}
\setcounter{enumi}{3}
\tightlist
\item
  In part (d) of the proof we assume that \(\mathbf{K} = \mathbf{R}\) . By § 8, no. 2, Theorem 2, \(\zeta(\mathrm{Aut}_1\mathrm{G})\) is a Lie subgroup of P. We transport the real Lie group structure on \(\zeta(\mathrm{Aut}_1\mathrm{G})\) to \(\mathrm{Aut}_1\mathrm{G}\) by means of \(\zeta^{-1}\) . Thus \(\mathrm{Aut}_1\mathrm{G}\) becomes a finite-dimensional Lie group.
\end{enumerate}

Let M be an analytic manifold, f a mapping of M into \(Aut_{1}G\) and \(\phi\) the mapping \((m,g)\mapsto f(m)g\) of \(M\times G\) into G. We have the following equivalences:

\(f\) analytic \(\Leftrightarrow\) the mappings \(m\mapsto (f(m))_i\) , where \(0\leqslant i\leqslant n\) , are analytic \(\Leftrightarrow \left\{ \begin{array}{ll}\text{the mappings } m\mapsto f(m)x_i\text{ of M into G, for } 1\leqslant i\leqslant n,\text{ are analytic} \\ \text{and} \\ \text{the mapping } (m,g)\mapsto f(m)g\text{ of M} \times \mathrm{G}_0\text{ into G is analytic} \end{array} \right.\) \(\Leftrightarrow \phi\) is analytic.

For \(w \in \mathrm{Aut}_1\mathrm{G}\) , \(\mathrm{L}(w) = \mathrm{L}(w_0)\) and hence the morphism \(w \mapsto \mathrm{L}(w)\) of \(\mathrm{Aut}_1\mathrm{G}\) into \(\mathrm{Aut}\, \mathrm{L}(\mathrm{G})\) is analytic. We see as in (b) that the law of infinitesimal operation associated with the law of operation of \(\mathrm{Aut}_1\mathrm{G}\) on \(\mathrm{G}\) is an isomorphism of \(\mathrm{L}(\mathrm{Aut}_1\mathrm{G})\) onto \(\mathfrak{a}\) .

Let C be a compact subset of \(G_{0}\) generating \(G_{0}\) . For a filter \(\Phi\) to converge to \(Id_{G}\) on \(Aut_{1}G\) , it is necessary and sufficient that \(\Phi|(C \cup \{x_{1}\} \cup \cdots \cup \{x_{n}\})\) and \(\Phi^{-1}|(C \cup \{x_{1}\} \cup \cdots \cup \{x_{n}\})\) converge uniformly to

\[
\operatorname{Id} _ {\mathbb {G}} \left| \left(\mathrm{C} \cup \{x _ {1} \} \cup \dots \cup \{x _ {n} \}\right). \right.
\]

The topology of \(\mathrm{Aut}_1\mathrm{G}\) is therefore the topology \(\mathcal{T}_{\beta}\) .

Clearly \(\mathrm{Aut}_1\mathrm{G}\) is open in Aut G with the topology \(\mathcal{T}_{\beta}\) . There exists on Aut G a Lie group structure compatible with this topology and inducing on \(\mathrm{Aut}_1\mathrm{G}\) the structure constructed above (§ 8, no. 1, Corollary 2 to Theorem 1). The fact that the Lie group Aut G has the properties of the theorem follows from the corresponding properties for \(\mathrm{Aut}_1\mathrm{G}\) .

\begin{enumerate}
\def\labelenumi{(\alph{enumi})}
\setcounter{enumi}{4}
\tightlist
\item
  In part (e) of the proof we assume that \(\mathbf{K} = \mathbf{C}\) . By (c) and Theorem 2 of §8, no. 2, there exists on \(\mathrm{Aut}_1\mathrm{G}\) a real Lie group structure such that \(\zeta\) is an isomorphism of \(\mathrm{Aut}_1\mathrm{G}\) onto a real Lie subgroup of \(\mathbb{P}\) .
\end{enumerate}

The law of operation \((w, g) \mapsto wg\) of \((\mathrm{Aut}_1\mathrm{G}) \times \mathrm{G}\) on \(\mathbf{G}\) is real analytic. Let \(\mathbf{D}\) be the associated law of infinitesimal operation. By Propositions 1 and 2 of no. 1, \(\mathrm{D}(\mathrm{L}(\mathrm{Aut}_1\mathrm{G})) = \mathfrak{a}\) .

For all \(\alpha \in \mathfrak{a}\) , let \(\alpha_0\) denote the restriction of \(\alpha\) to \(\mathrm{G}_0\) ; it is an infinitesimal automorphism of \(\mathrm{G}_0\) which we identify, because of part (b) of the proof, with an element of \(\mathrm{L}(\mathrm{Aut}\mathrm{G}_0)\) . For \(1 \leqslant i \leqslant n\) , we write

\[
\alpha_ {i} = x _ {i} ^ {- 1} \alpha (x _ {i}) \in \mathrm{L} (\mathrm{G}) = \mathrm{L} (\mathrm{G} _ {0}).
\]

Finally, we write \(f(\alpha) = ((\alpha_1, \ldots, \alpha_n), \alpha_0) \in \mathrm{L}(\mathrm{P})\) . Then \(f\) is a \(\mathbf{C}\) -linear mapping of \(\mathfrak{a}\) into \(\mathrm{L}(\mathrm{P})\) .

On the other hand, clearly \(\mathrm{L}(\zeta) = f\circ \mathrm{D}\) . Hence \(\mathrm{L}(\zeta)(\mathrm{L}(\mathrm{Aut}_1\mathrm{G})) = f(a)\) is a complex vector subspace of \(\mathrm{L}(P)\) . By Proposition 2 of § 4, no. 2, \(\zeta (\mathrm{Aut}_1\mathrm{G})\) is a complex Lie subgroup of \(P\) and we can proceed exactly as in (d): we transport the complex Lie group structure on \(\zeta (\mathrm{Aut}_1\mathrm{G})\) to \(\mathrm{Aut}_1\mathrm{G}\) by means of \(\zeta^{-1}\) and we see as in (d) that \(\mathrm{Aut}_1\mathrm{G}\) has the properties analogous to properties (i), (ii), (iii), (v) and (vi) of the theorem.

Clearly \(\mathrm{Aut}_1\mathrm{G}\) is open in Aut G with the topology \(\mathcal{T}_{\beta}\) . Let \(w\in \mathrm{Aut}G\) . Let \(\sigma\) be the automorphism \(v\mapsto wvw^{-1}\) of \(\mathrm{Aut}_1\mathrm{G}\) . It is real analytic (§ 8, no. 1, Theorem 1), \(\mathrm{L}(\sigma)\) is an \(\mathbf{R}\) -automorphism of \(\mathrm{L}(\mathrm{Aut}_1\mathrm{G})\) and

\[
\mathrm{D} \circ \mathrm{L} (\mathrm{Aut} _ {1} \mathrm{G}) \circ \mathrm{D} ^ {- 1}
\]

is an R-automorphism of a. This automorphism is also the automorphism of a derived from w by transport of structure; as w is K-analytic, we see that \(L(\sigma)\) is K-linear. Hence \(\sigma\) is K-analytic ( \(\S\) 3, no. 8, Proposition 32). By \(\S\) 1, no. 9, Proposition 18, there exists on Aut G one and only one Lie K-group structure such that \(Aut_1G\) is an open Lie subgroup of Aut G. The fact that this structure has the properties of the theorem follows from the corresponding properties for \(Aut_1G\) .

COROLLARY 1. Let G be a finite-dimensional real Lie group and \(G_{0}\) its identity component. Suppose that G is generated by \(G_{0}\) and a finite number of elements. Then Aut G has the topology \(T_{\beta}\) and is a finite-dimensional real Lie group.

COROLLARY 2. Let G be a semi-simple connected real or complex Lie group. The group Int G is the identity component of Aut G.

The mapping \(u \mapsto \mathrm{L}(u)\) is an isomorphism of Aut G onto a Lie subgroup of Aut L(G) (Theorem 1). The image of Int G under this isomorphism is Ad G. But Ad G is the identity component of Aut L(G) (§ 9, no. 8, Proposition 30 (ii)).

